\section{Method}
\label{sec:method}

We treat each $F$-frame GOP as the unit of compression. The decoder
reconstructs the GOP by running a pretrained rectified-flow video
generator forward for $T$ discrete steps; the encoder selects per-step
codebook indices so that the resulting trajectory matches the ground-truth
GOP. Section~\ref{sec:method-pipeline} states this codec backbone in
codec-centric terms (inherited from~\cite{zeng2026gvcc,vaisman2025turboddcm}
in compact form), and the remaining subsections describe the two
acceleration mechanisms.

\subsection{Codec backbone}
\label{sec:method-pipeline}

\paragraph{Encoder/decoder pipeline.}
A pretrained 3D VAE maps a GOP $\mathbf{V}\!\in\!\mathbb{R}^{F\times H\times W\times 3}$
to a latent $\mathbf{x}_0\!\in\!\mathbb{R}^{F_l\times H_l\times W_l\times C}$.
The decoder is a rectified-flow video generator whose ODE
$\dot{\mathbf{x}}_t \!=\! \mathbf{u}_\theta(\mathbf{x}_t, t)$ transports
samples from $t\!=\!1$ (pure noise) to $t\!=\!0$ (clean latent). We
discretize this trajectory at $T$ timesteps $\{t_k\}_{k=1}^T$ and inject
stochastic innovations at the first $T_\text{sde}$ of them, leaving an
$N\!=\!T\!-\!T_\text{sde}$-step deterministic ODE tail:
\begin{equation}
\mathbf{x}_{t_{k+1}} = \mathbf{x}_{t_k} \!-\! f_{t_k}\!\Delta t_k + g_{t_k}\sqrt{\Delta t_k}\,\mathbf{z}^*_{k},
\label{eq:sde-step}
\end{equation}
where $f_{t_k}, g_{t_k}$ are the SDE drift and diffusion coefficients
derived from $\mathbf{u}_\theta$ as in~\cite{zeng2026gvcc}, and
$\mathbf{z}^*_{k}$ is the per-step innovation that the encoder
\emph{transmits} via a $\log_2(K)$-bit codebook index together with a sign
bit, per latent frame. The full bitstream therefore consists of
$T_\text{sde}\!\cdot\! F_l \!\cdot\! M\!\cdot\!(\log_2 K + 1)$ codebook
bits per GOP, plus any boundary-frame bits required by the conditioning
mode (Section~\ref{sec:method-gop}).

\paragraph{Codebook construction.}
Per-step innovations are selected via the multi-atom construction of
Turbo-DDCM~\cite{vaisman2025turboddcm}: at step $k$ and latent frame
$f$ the encoder generates $K$ i.i.d.\ Gaussian atoms from a shared
deterministic seed, transmits the $M$ atoms with the largest absolute
inner product against the per-frame residual
$\mathbf{r}_{k,f}\!=\!\mathbf{x}_{0,f}\!-\!\hat{\mathbf{x}}_{0|t_k,f}$
(together with their signs), and reconstructs $\mathbf{z}^*_{k,f}$ as
the corresponding unit-variance sum on the decoder side. The cost is
dominated by $T$ video-generator forwards per GOP---approximately
$700$~s at $1080$p on the 14B Wan~2.1 model. The two mechanisms below
reduce this cost without changing the codec backbone.

\subsection{Trajectory-preserving acceleration principle}
\label{sec:method-principle}

The key distinction from ordinary fast sampling is that GVCC is a
\emph{codec}: the transmitted bitstream specifies the stochastic
innovations $\mathbf{z}^*_k$ at every SDE step, and the decoder must
reproduce the same trajectory as the encoder. Directly reducing the
number of SDE steps $T_\text{sde}$ would therefore remove not only solver
work but also the opportunities to apply the transmitted codebook
correction; the codec channel itself would shrink.

GVCCTurbo instead keeps the codebook correction grid intact and reduces
the per-step cost on a subset of steps in two ways. \emph{Spatially},
early high-noise steps may run on a smaller latent grid, provided the
stage target remains a projection of the final high-resolution latent so
that the lifted trajectory is still a trajectory toward the same clean
latent $\mathbf{x}_0$ (Section~\ref{sec:method-mr}). \emph{Temporally},
selected per-step generator evaluations may be replaced by a
deterministic prediction derived from a decoder-reproducible cache, with
the SDE codebook correction (Eq.~\ref{eq:sde-step}) still applied at
every skipped step (Section~\ref{sec:method-skip}). The rest of this
section instantiates these two reductions; the common constraint is that
neither changes the codebook correction channel that carries the
compressed bits.

\subsection{Multi-resolution coding with a unified target}
\label{sec:method-mr}

\paragraph{Resolution schedule.}
We replace the single-resolution trajectory with an $S$-stage progressive
schedule, $s\!=\!0,\ldots,S\!-\!1$, where stage $s$ operates on a latent of
spatial size $(H_l^{(s)}, W_l^{(s)})$ with
$H_l^{(0)} \!\le\! \cdots \!\le\! H_l^{(S-1)} \!=\! H_l$. Stage $s$
denoises for $\tau_s$ codebook steps; at transition step $\tau_{s\to s+1}$
the trajectory is lifted to the next resolution via a 2D DCT zero-pad in
the new band:
\begin{equation}
\bar{\mathbf{x}}^{(s+1)}_{t_k} = \mathcal{T}^{-1}\!\Bigl(
\text{Embed}\bigl(\mathcal{T}(\mathbf{x}^{(s)}_{t_k})\bigr) + t_k\,\boldsymbol{\epsilon}^{(s+1)}_H
\Bigr),
\label{eq:transition}
\end{equation}
where $\mathcal{T}$ is the spatial DCT, $\text{Embed}$ pads the new
high-frequency band with zeros, and $\boldsymbol{\epsilon}^{(s+1)}_H$ is a
fresh Gaussian sample on that band. Equation~\ref{eq:transition} is
spectrally correct: the low band passes through unchanged and the new band
matches the marginal distribution expected by the generator at the new
resolution and timestep $t_k$.

\paragraph{The target-domain mismatch.}
For the encoder to drive the low-stage trajectory, it needs a low-stage
ground-truth latent $\mathbf{x}_0^{(s)}$. The natural choice is
\begin{equation}
\mathbf{x}_0^{(s),\,\text{naive}} \;=\; \mathrm{VAE}^{(s)}\bigl(\mathrm{down}(\mathbf{V})\bigr),
\label{eq:naive-target}
\end{equation}
i.e., spatially downsample the pixel video and re-encode at the lower
resolution---what a single-resolution codec would do at each scale
independently. The decoder, however, does \emph{not} produce this object
at the end of stage $s$: it produces the trajectory whose low band, after
Eq.~\ref{eq:transition} lifts it, agrees with the DCT-low projection of
the \emph{high-resolution} latent,
$\mathcal{T}_L(\mathrm{VAE}(\mathbf{V}))$.

The two are different objects:
$\mathrm{VAE}^{(s)}(\mathrm{down}(\mathbf{V})) \neq \mathcal{T}_L(\mathrm{VAE}(\mathbf{V}))$,
because the VAE encoder is nonlinear and its receptive field on the
downsampled pixel support is not the corresponding DCT projection of its
receptive field on the full-resolution support. On UVG with the Wan~2.1
VAE the relative gap is approximately $0.41$ at $t\!=\!0$ on the low
band. The encoder under Eq.~\ref{eq:naive-target} therefore drives the
trajectory toward an object the decoder cannot produce; whatever
high-band content is added at the transition cannot recover the
mismatch.

\paragraph{Fix: unify both stages in the high-resolution latent domain.}
We replace Eq.~\ref{eq:naive-target} with
\begin{equation}
\mathbf{x}_0^{(s),\,\text{ours}} \;=\; \mathcal{T}_L\bigl(\mathrm{VAE}(\mathbf{V})\bigr),
\label{eq:lp-target}
\end{equation}
i.e., encode the pixel video \emph{once} at full resolution and project to
the low band by DCT truncation. Both stages now reference the same latent
$\mathbf{x}_0 \!=\! \mathrm{VAE}(\mathbf{V})$, observed at different spectral
bandwidths. Transition (Eq.~\ref{eq:transition}) stops being a domain
change and becomes pure band expansion: stage $s$ converges to
$\mathcal{T}_L(\mathbf{x}_0)$, and stage $s\!+\!1$ adds
$\mathcal{T}_H(\mathbf{x}_0)$, recovering the same $\mathbf{x}_0$ that the
encoder is matching against. The bitstream is unchanged; only the
encoder's matching target moves to a spectrally consistent reference.

\paragraph{Transition codebook.}
At the transition step $\tau_{s\to s+1}$, the newly opened high band is
by default filled with pure noise $t_k\boldsymbol{\epsilon}_H$
(Eq.~\ref{eq:transition}). The encoder knows the target high-band
content---it is $(1-t_k)$ times the DCT-high projection of
$\mathbf{x}_0$---and we transmit one additional codebook atom per latent
frame at this step, selecting atoms by their inner product with this
known residual. The Gaussian noise floor is left untouched; only the
clean-target component carries bits. This costs $<1\%$ of the per-GOP
bit budget and recovers up to $0.4$~dB PSNR over random high-band fill
(Section~\ref{sec:ablations}).

\subsection{Predictive skip via clean-endpoint caching}
\label{sec:method-skip}

The dominant per-step cost is the video-generator evaluation
$\mathbf{u}_\theta(\mathbf{x}_t, t)$. We aim to replace approximately
half of these evaluations along the codebook-driven trajectory.

\paragraph{Why velocity caching fails.}
The naive choice is to cache the velocity itself and reuse the previous
step's direction
$\mathbf{u}^{\,\text{naive}}_{k+1} = \mathbf{u}_\theta(\mathbf{x}_{t_k},t_k)$.
This silently drifts: $\mathbf{u}_\theta$ is a function of both
$\mathbf{x}$ and $t$, and the trajectory has moved away from
$\mathbf{x}_{t_k}$ by both a deterministic drift and the encoder-injected
codebook innovation between steps $k$ and $k\!+\!1$ in
Eq.~\ref{eq:sde-step}. The cached direction is computed at a stale state,
not the current one.

\paragraph{The fix: cache the trajectory endpoint.}
Under the linear flow interpolant
$\mathbf{x}_t = (1-t)\,\mathbf{x}_0 + t\,\boldsymbol{\epsilon}$, the
generator's MMSE estimate of $\mathbf{x}_0$ at any point on the
trajectory is
\begin{equation}
\hat{\mathbf{x}}_{0|t_k} \;=\; \mathbf{x}_{t_k} - t_k\,\mathbf{u}_\theta(\mathbf{x}_{t_k}, t_k).
\label{eq:x0-hat}
\end{equation}
$\hat{\mathbf{x}}_{0|t_k}$ is, geometrically, where the model thinks the
trajectory ends. We cache this endpoint instead. When we skip the
generator call at step $k\!+\!1$, we recover an approximate velocity from
the cached endpoint and the \emph{current} state and time:
\begin{equation}
\tilde{\mathbf{u}}_{k+1} \;=\; \bigl(\mathbf{x}_{t_{k+1}} - \hat{\mathbf{x}}_{0|t_k}\bigr) \,/\, t_{k+1}.
\label{eq:u-tilde}
\end{equation}
This is valid up to a local linearization around $\hat{\mathbf{x}}_{0|t_k}$:
because the clean endpoint is invariant along the deterministic part of
the trajectory under linear flow, only the codebook innovation between
steps $k$ and $k\!+\!1$ moves the endpoint, and that movement is small
relative to the cached value. Crucially, the full SDE/codebook correction
in Eq.~\ref{eq:sde-step} is preserved at the skipped step; only the
generator query is replaced.

\paragraph{Schedule and warmup.}
We adopt the fixed alternating schedule
$k\;\text{is run}\Leftrightarrow k\!\equiv\!k_0\pmod 2$, where $k_0$ is the
first SDE step in the current resolution stage. At every transition the
cache is invalidated (the spatial resolution and therefore the operating
domain changes), and we additionally delay the first post-transition skip
by one step (\emph{transition warmup}): the cache is rebuilt at step
$\tau_{s\to s+1}\!+\!1$ before being reused at step
$\tau_{s\to s+1}\!+\!2$. The number of generator calls is unchanged
relative to the no-warmup schedule; the only difference is which step's
output drives the first post-transition skip.

\paragraph{One-step radius.}
Empirically the cache is locally valid for one skipped step. We tested
adaptive variants that allowed two or more consecutive skips when a
running per-step error estimate fell below a threshold; quality collapsed
as soon as the cap of two was actually triggered (Section~\ref{sec:ablations}).
The mechanism is recursive cache staleness: the second skipped step plugs
an already-extrapolated state back into the same linearization, and the
error compounds. This radius-of-one observation rules out a large family
of more aggressive schedules and identifies the fixed alternating pattern
as Pareto-optimal under this caching scheme.

\subsection{GOP-level conditioning modes}
\label{sec:method-gop}

We retain the three GOP-conditioning modes of GVCC~\cite{zeng2026gvcc}
unchanged. \textbf{T2V} generates each GOP unconditionally from text;
the bitstream is codebook-only and the operating point is the lowest in
our set. \textbf{I2V} takes the ground-truth first frame of the sequence
as a free anchor and chains subsequent GOPs autoregressively using the
previous GOP's decoded last frame, plus a small quantized tail-residual
correction (8-bit, last latent frame) to suppress drift.
\textbf{FLF2V} conditions each GOP on both its endpoint frames,
compressed with CompressAI~\cite{begaint2020compressai} at quality 4;
boundaries are shared across consecutive GOPs ($N\!+\!1$ frames for $N$
GOPs), amortizing the cost over long sequences. The three modes occupy
distinct points on the rate-quality plane; Section~\ref{sec:experiments}
evaluates them.
